\chapter{Electronics Design}
\label{ch:electronics}

\section{RF PCB Design}
\label{sec:rf_pcb}

The RF PCB implements the front-end architecture responsible for signal transmission, reception, and communication with the host system. This section describes the electrical design of the PCB, including system architecture, component selection, and interconnections between subsystems.

\subsection{System Overview}
\label{subsec:rf_system_overview}

The RF PCB is based on a microcontroller-controlled dual-transceiver architecture supporting both the VHF (144\,MHz) and UHF (433\,MHz) frequency bands.

The design consists of the following functional blocks:

\begin{itemize}
\item A microcontroller for system control and data handling
\item Two RF transceivers (VHF and UHF)
\item Dedicated SPI communication interfaces
\item A UART interface to the Raspberry Pi host
\item Power distribution and filtering circuitry
\end{itemize}

Figure~X illustrates the high-level architecture of the RF PCB and the interaction between these subsystems.

\subsection{Microcontroller Subsystem}
\label{subsec:rf_microcontroller}

The RF PCB is controlled by a Raspberry Pi Pico microcontroller based on the RP2040. This device acts as the central control unit of the system.

Its primary responsibilities include:

\begin{itemize}
\item Configuration of both RF transceivers via SPI
\item Handling of transmit and receive data
\item Communication with the host system via UART
\item Control of auxiliary signals such as chip select and interrupt lines
\end{itemize}

The RP2040 was selected due to its availability of multiple SPI peripherals, allowing independent communication with each RF transceiver. Additionally, its 3.3,V logic level ensures direct compatibility with the RF components without requiring level shifting.

Furthermore, the platform supports rapid firmware development and iteration through a well-established software ecosystem, enabling efficient testing and debugging during the development process.

The GPIO assignment of the RP2040 was selected to utilize two independent SPI peripherals while maintaining a simple and deterministic interface to the host system.

\begin{table}[H]
\centering
\caption{RP2040 Pin Assignment for RF PCB}
\begin{tabular}{|c|c|c|}
\hline
\textbf{GPIO} & \textbf{Function} & \textbf{Connected To} \\
\hline
GPIO0  & UART TX & Raspberry Pi RX (J3) \\
GPIO1  & UART RX & Raspberry Pi TX (J3) \\
\hline
GPIO10 & SPI (UHF) MOSI & CC1200 UHF \\
GPIO11 & SPI (UHF) MISO & CC1200 UHF \\
GPIO12 & SPI (UHF) SCLK & CC1200 UHF \\
GPIO13 & SPI (UHF) CS   & CC1200 UHF \\
\hline
GPIO18 & SPI (VHF) MOSI & CC1200 VHF \\
GPIO19 & SPI (VHF) MISO & CC1200 VHF \\
GPIO16 & SPI (VHF) SCLK & CC1200 VHF \\
GPIO17 & SPI (VHF) CS   & CC1200 VHF \\
\hline
\end{tabular}
\label{tab:rp2040_pinout}
\end{table}

\subsection{RF Transceiver Subsystem}
\label{subsec:rf_transceivers}

The RF PCB employs two CC1200 transceivers to enable dual-band operation:

\begin{itemize}
\item One transceiver configured for the 144\,MHz VHF band
\item One transceiver configured for the 433\,MHz UHF band
\end{itemize}

Each transceiver operates independently and includes its own supporting circuitry:

\begin{itemize}
\item A 40,MHz crystal oscillator for frequency reference
\item External matching and filtering components
\item Decoupling capacitors for supply stabilization
\end{itemize}

The use of two dedicated transceivers eliminates the need for RF switching networks, simplifying the signal path and avoiding additional insertion losses.

Key control and interface signals include SPI communication lines, interrupt outputs (GDO pins), and device control signals such as reset.

This independent operation allows simultaneous configuration and rapid switching between frequency bands at the system level.

\subsection{Communication Interfaces}
\label{subsec:rf_interfaces}

\subsubsection{SPI Communication}
\label{subsubsec:rf_spi}

Each RF transceiver is connected to the microcontroller via a dedicated hardware SPI interface on the RP2040, allowing independent communication with both devices.

The SPI interfaces consist of the following signals:

\begin{itemize}
\item Serial clock (SCLK)
\item Master-out slave-in (MOSI)
\item Master-in slave-out (MISO)
\item Chip select (CS)
\end{itemize}

Interrupt lines from each transceiver are connected to the microcontroller to signal events such as packet reception or transmission completion.

Using separate SPI buses avoids contention and ensures deterministic communication timing.

\subsubsection{UART Interface to Host}
\label{subsubsec:rf_uart}

Communication between the RF PCB and the Raspberry Pi host is implemented via a UART interface.

This interface is used for:

\begin{itemize}
\item Transferring received RF data to the host
\item Receiving configuration and control commands
\item Debugging and diagnostic output
\end{itemize}

The UART operates at 3.3\,V logic levels and provides a simple and robust communication channel.


\subsection{Power Distribution Network}
\label{subsec:rf_power}

The RF PCB is powered from a 24\,V input supplied via an XT30 connector. This input voltage is converted to a regulated 5\,V rail using a buck converter. The resulting 5\,V supply is used both to power the Raspberry Pi host through the board header and to supply the VSYS input of the Raspberry Pi Pico microcontroller.

Using a single 5\,V rail for both the host and the onboard controller simplifies the overall power architecture and reduces the required component count while maintaining sufficient efficiency.

The Raspberry Pi Pico generates the required 3.3\,V rail using its onboard power supply circuitry, which is based on a switching regulator. This 3.3\,V supply powers both CC1200 transceivers as well as the associated digital circuitry on the RF PCB. While switching regulators typically introduce more noise compared to dedicated low-noise linear regulators, the integrated solution provides adequate current capability and efficiency for the system.

Given the noise sensitivity of RF circuits, particular attention was paid to supply integrity. Local decoupling capacitors were placed close to the supply pins of each active component to reduce high-frequency noise, provide transient current during switching events, and improve overall system stability. In addition, careful PCB layout practices were applied to minimize coupling of digital switching noise into the RF signal paths.

\subsubsection{Power Supply Noise Evaluation}
\label{subsubsec:rf_power_measurement}

Oscilloscope measurements were performed at the 3.3\,V pins of the CC1200 transceivers to assess voltage stability and switching noise characteristics. The measurements were conducted under multiple operating conditions, including idle, receive, and transmit modes.

The observed DC supply voltage was approximately \textbf{[X.XX\,V]}. The measured peak-to-peak ripple on the 3.3\,V rail was approximately \textbf{[XX\,mV\textsubscript{pp}]} under nominal operation, with transient deviations up to \textbf{[XX\,mV]} observed during mode transitions.

No significant voltage drops or disturbances were observed that could affect the stability of the RF transceivers. Furthermore, the system demonstrated reliable communication during operation, indicating that the supply quality was sufficient for the intended application.

These results suggest that, despite the use of a switching regulator, the implemented power distribution network, in combination with local decoupling, provides adequate noise suppression for this prototype. For future iterations, the use of a dedicated low-noise LDO for the RF section may be considered to further improve supply integrity and RF performance.

Figure~X illustrates the power distribution architecture of the RF PCB, including the conversion from the 24\,V input to the regulated 5\,V and 3.3\,V supply rails.

\subsection{Summary}
\label{subsec:rf_summary}

The RF PCB electrical design integrates a microcontroller, dual RF transceivers, and dedicated communication interfaces into a modular and robust architecture. This design forms the foundation for the RF layout considerations discussed in the following chapter.

\section{Azimuth/Elevation (AZ/EL) PCB Design}
\label{sec:azel_pcb}

This section will describe the electrical design of the azimuth and elevation control PCB, including motor control, sensing, communication interfaces, and power distribution.

% (To be completed)
